\chapter{Continuously differentiable pieces of Lipschitz images}\label{sec:rectifiable-c1-decomposition}
The tangent theorem for a rectifiable set begins with a decomposition into smooth pieces, rather than with a separation assumption between those pieces. We expand \cite{Sim18}, Chapter 2, Theorems 1.5--1.6, printed pp.~46--48, and Chapter 3, Lemma 1.2 and Remark 1.3, pp.~85--86. The Whitney extension theorem and Euclidean area formula remain explicit external inputs. The approximation concerns agreement outside a set of small measure, not uniform approximation of derivatives everywhere.

\begin{definition}[A compatible first-order jet on a closed set]\label{def:whitney-first-jet}
Let $C\subset\mathbb R^m$ be closed. A scalar first-order jet consists of continuous maps $h:C\to\mathbb R$ and $a:C\to\mathbb R^m$. It is compatible if for each compact $K\subset C$,
\[
 \lim_{\delta\downarrow0}\sup_{\substack{x,y\in K\\0<|x-y|<\delta}}
 \frac{|h(y)-h(x)-a(x)\cdot(y-x)|}{|y-x|}=0.
\]
The supremum over an empty set is zero. Uniformity over pairs in a compact set is part of the condition. This is the condition in \cite{Sim18}, Chapter 2, Theorem 1.6.
\end{definition}

\begin{theorem}[Whitney extension for first-order jets]\label{thm:whitney-c1-extension}
Every compatible jet in Definition \ref{def:whitney-first-jet} is the restriction of a function $g\in C^1(\mathbb R^m)$: $g=h$ and $Dg=a$ on $C$. Conversely, the restriction of a $C^1$ function gives a compatible jet. The existence assertion is a cited theorem, \cite{Sim18}, Chapter 2, Theorem 1.6, pp.~46--47, where the proof is referred to \cite{Fed69}. Its extension construction is not proved here. For necessity, the segment formula for $g(y)-g(x)$ and uniform continuity of $Dg$ on the compact convex hull of $K$ give the displayed remainder estimate.
\end{theorem}

\begin{lemma}[Measurable differentiability defects]\label{lem:measurable-differentiability-defect}
Let $f:\mathbb R^m\to\mathbb R$ be Lipschitz and let $a(x)$ be its derivative at differentiability points, extended by zero on the exceptional set. It can be chosen Borel. Put
\[
 d_k(x)=\sup_{0<|z|<1/k}
 \frac{|f(x+z)-f(x)-a(x)\cdot z|}{|z|}.
\]
Then $d_k$ is Borel, is bounded by twice the Lipschitz constant at differentiability points, and decreases to zero at every such point. This is the defect used in the proof of \cite{Sim18}, Theorem 1.5, p.~48.
\end{lemma}
\begin{proof}
Coordinate difference quotients are continuous functions of $x$, so the set where their limits exist is Borel and their limits there are Borel. Rademacher's theorem, Theorem \ref{thm:rademacher}, supplies a full-measure differentiability set. The full differentiability condition can itself be expressed by the vanishing of the defects below using the coordinate limits, hence is Borel. Set $a=0$ off that set. At a fixed $x$ the quotient is continuous in $z\ne0$, so its supremum can be taken over nonzero rational vectors with length less than $1/k$. It is consequently a countable supremum of Borel functions. At a differentiability point $|a|\leq\operatorname{Lip}f$ and the defining differentiability limit proves the last two assertions. The shrinking domains of the suprema give monotonicity in $k$.
\end{proof}

\begin{lemma}[Compact subsets with continuous and uniform jets]\label{lem:compact-good-lipschitz-jets}
For a Lipschitz scalar $f$, a bounded Lebesgue-measurable set $A\subset\mathbb R^m$, and $\varepsilon>0$, there is a compact $C\subset A$ consisting of differentiability points with $|A\setminus C|<\varepsilon$, such that $Df|C$ is continuous and the defects in Lemma \ref{lem:measurable-differentiability-defect} tend uniformly to zero on $C$. The restricted jet $(f,Df)|C$ is compatible.
\end{lemma}
\begin{proof}
Discard the null nondifferentiability set. Egoroff's theorem on the finite-measure set $A$ gives a measurable subset with complement of measure less than $\varepsilon/2$ on which $d_k\to0$ uniformly. Lusin's theorem and inner regularity give a compact subset of it, with additional loss less than $\varepsilon/2$, on which $Df$ is continuous. These two elementary finite-measure theorems are the background inputs also used in \cite{Sim18}, p.~48. Continuity of $f$ is automatic. For $x,y\in C$ with $0<|x-y|<1/k$, the normalized jet remainder is at most $d_k(x)$. Uniform convergence therefore gives the compatibility condition. This step needs neither an extension of the jet nor an area formula.
\end{proof}

\begin{theorem}[Lipschitz functions agree with a smooth function off a small set]\label{thm:lipschitz-c1-lusin-approximation}
If $f:\mathbb R^m\to\mathbb R$ is Lipschitz and $\varepsilon>0$, there exists $g\in C^1(\mathbb R^m)$ such that
\[
 \mathcal L^m\bigl(\{f\ne g\}\cup\{Df\ne Dg\}\bigr)<\varepsilon,
\]
where the derivative comparison ignores the null nondifferentiability set of $f$. This is \cite{Sim18}, Chapter 2, Theorem 1.5, with the Whitney extension step isolated above.
\end{theorem}
\begin{proof}
Partition space into the bounded shells $A_j=\{j-1\leq|x|<j\}$, $j\geq1$. Apply Lemma \ref{lem:compact-good-lipschitz-jets} to obtain compact $C_j\subset A_j$ with loss less than $\varepsilon 2^{-j}$, continuous restricted derivative, and uniform differentiability defects. The family $\{C_j\}$ is locally finite: a bounded neighborhood meets only finitely many shells. Hence $C=\bigcup_jC_j$ is closed. The compact sets are disjoint. In any bounded neighborhood there are only finitely many of them, so their pairwise distances are positive whenever both are nonempty. Consequently $Df|C$ is continuous locally and the defects tend uniformly to zero on each compact subset of $C$, being the maximum of finitely many uniform defects. The jet is compatible.

Theorem \ref{thm:whitney-c1-extension} supplies $g\in C^1$ with $g=f$ and $Dg=Df$ on $C$. The union of the two disagreement sets is contained, up to the stated null set, in $\mathbb R^m\setminus C$. Its measure is less than $\sum_j\varepsilon2^{-j}=\varepsilon$. No finiteness of the total ambient measure or compact support of $f$ is assumed.
\end{proof}

\begin{corollary}[Vector-valued agreement on one common good set]\label{cor:vector-c1-lusin-approximation}
If $f:\mathbb R^m\to\mathbb R^n$ is Lipschitz and $\varepsilon>0$, there is $g\in C^1(\mathbb R^m;\mathbb R^n)$ such that $f=g$ and $Df=Dg$ outside a set of Lebesgue measure less than $\varepsilon$.
\end{corollary}
\begin{proof}
Apply Theorem \ref{thm:lipschitz-c1-lusin-approximation} to each of the $n$ scalar components, with exceptional measure less than $\varepsilon/n$. Collect the resulting functions into $g$. Outside the union of their exceptional sets every component and every derivative agrees. Subadditivity gives the claimed measure bound. The common set is essential when using the rank of the vector-valued derivative.
\end{proof}

\begin{theorem}[Countable decomposition into embedded smooth manifolds]\label{thm:rectifiable-c1-decomposition}
For $1\leq m\leq n$, an $\mathcal H^m$-measurable $M\subset\mathbb R^n$ is countably $m$-rectifiable if and only if it is contained, apart from an $\mathcal H^m$-null set, in a countable union of embedded $m$-dimensional $C^1$ submanifolds. No local finiteness of $\mathcal H^m\llcorner M$ is required for this equivalence. This is \cite{Sim18}, Chapter 3, Lemma 1.2, pp.~85--86.
\end{theorem}
\begin{proof}
First take a Lipschitz parametrization $f:\mathbb R^m\to\mathbb R^n$. For each positive integer $i$, Corollary \ref{cor:vector-c1-lusin-approximation} gives a $C^1$ map $g_i$ agreeing with $f$ outside a set of measure less than $1/i$. The set $Z$ on which none of the equalities $f=g_i$ holds has measure zero, since it is contained in each disagreement set. The elementary Hausdorff-measure Lipschitz bound gives $\mathcal H^m(f(Z))=0$. Thus $f(\mathbb R^m)$ is covered, apart from that null image, by the images $g_i(\mathbb R^m)$.

For a fixed $g_i$, its critical set $C_i=\{\operatorname{rank}Dg_i<m\}$ has $\mathcal H^m(g_i(C_i))=0$ by the cited area formula, Theorem \ref{thm:area-formula}: its Jacobian vanishes, while the multiplicity is at least one on its image. To apply the Lipschitz version of that formula, exhaust the domain by bounded balls on whose closures $Dg_i$ is bounded, apply it on each ball, and take the countable union of the null critical images. At every point outside $C_i$, some $m$ coordinate functions have invertible derivative. The inverse function theorem writes the image of a sufficiently small neighborhood as an embedded $C^1$ graph over those coordinates; injectivity follows from that same coordinate inverse. A countable subcover of the open full-rank domain gives countably many such graphs. Apply this argument to every Lipschitz map in Definition \ref{def:rectifiable-set}, and collect the countably many null sets and charts.

Conversely each embedded $C^1$ manifold has a countable cover by chart images with bounded derivatives on relatively compact base balls. Each graph chart is Lipschitz there and extends coordinatewise to a Lipschitz map on all of $\mathbb R^m$ by Proposition \ref{prop:mcshane-extension}. The vector extension may increase the constant by the square root of its target dimension, which does not affect countable rectifiability. Add the given exceptional null set. The proof uses Whitney extension, the inverse function theorem and the Euclidean area formula; it does not use the tangent-plane characterization it will help establish.
\end{proof}

\begin{lemma}[A countable cover by finite-area graph patches]\label{lem:finite-area-c1-patches}
An embedded $m$-dimensional $C^1$ submanifold of $\mathbb R^n$ has a countable cover by relatively open graph patches $N_j$ of finite $\mathcal H^m$ measure. Every point of each patch is an interior point of its parametrization domain; the closed base ball of each patch lies inside a larger chart with bounded derivative. For a measurable rectifiable set $M$, these patches can be made into a disjoint measurable partition of $M$, apart from a null set, by assigning a point to the first patch containing it.
\end{lemma}
\begin{proof}
Use local graph charts from the inverse function theorem and choose smaller open base balls with closures inside their original domains. A countable subcover exists by the countable basis of the ambient space. The derivative is bounded on each closed base ball. The graph Jacobian is therefore bounded, and the area formula gives finite measure because the base ball has finite volume. The graph of a continuous map over an open domain is Borel. Number the patches, intersect them with $M$, and remove all preceding pieces to obtain disjoint measurable sets. The null exceptional set from Theorem \ref{thm:rectifiable-c1-decomposition} remains null. Artificial patch edges are excluded by using open base balls; every covered point is interior to at least one patch. This avoids assuming that an entire nonproper submanifold has locally finite ambient area.
\end{proof}
